\documentclass{../tikz_figure}
\begin{document}
\begin{tikzpicture}

    % 设置波源位置
    \def\A{(-2,0)}
    \def\B{(2,0)}

    % 绘制波源点
    \filldraw[black] \A circle (2pt) node[above] {A};
    \filldraw[black] \B circle (2pt) node[above] {B};

    % 波前的最大半径
    \def\maxcount{7}    % 圆的数量
    \def\interval{0.5}   % 间距

    % 从波源A画同心圆
    \foreach \i in {1,...,\maxcount} {
            \pgfmathparse{\i*\interval}
            \let\radius\pgfmathresult
            \ifodd\i
                \draw[blue] \A circle (\radius); % 实线 - 波峰
            \else
                \draw[blue,dashed] \A circle (\radius); % 虚线 - 波谷
            \fi
        }

    % 从波源B画同心圆
    \foreach \i in {1,...,\maxcount} {
            \pgfmathparse{\i*\interval}
            \let\radius\pgfmathresult
            \ifodd\i
                \draw[red] \B circle (\radius); % 实线 - 波峰
            \else
                \draw[red,dashed] \B circle (\radius); % 虚线 - 波谷
            \fi
        }

    % 标注说明
    \node at (0,4) {\textbf{干涉图样}};
    \draw[dashed, black] (4,3.3) -- +(0.5,0) node[right]{波谷（虚线）};
    \draw[black] (4,2.9) -- +(0.5,0) node[right]{波峰（实线）};

\end{tikzpicture}
\end{document}